%==============================================================================
%  Effect of an authored joint-limit prior WITHOUT 3D supervision
%  SMILify / SMILySTICKS -- internship report  (2D-only study)
%
%  STATUS: PLACEHOLDER. Every number, figure and conclusion marked in red
%  (\PH{...}, \TBD, \placeholderfig) must be filled from
%  prior_study_results_2d/ once the three arms have finished training.
%  Set \finaltrue below to hide the red markers' colour for the final build.
%
%  Compile:  pdflatex main.tex  (twice, for cross-references)
%==============================================================================
\documentclass[11pt,a4paper]{article}

\usepackage[T1]{fontenc}
\usepackage[utf8]{inputenc}
\usepackage{lmodern}
\usepackage[english]{babel}
\usepackage[margin=2.5cm]{geometry}
\usepackage{amsmath,amssymb}
\usepackage{graphicx}
\usepackage{booktabs}
\usepackage{multirow}
\usepackage{siunitx}
\usepackage{caption}
\usepackage{subcaption}
\usepackage{enumitem}
\usepackage{xcolor}
\usepackage[hidelinks]{hyperref}
\usepackage{cleveref}

\sisetup{
  detect-weight = true,
  detect-family = true
}

\captionsetup{font=small,labelfont=bf}
\captionsetup[sub]{font=small}
\graphicspath{{figures_2d/}{figures/}}
\setlength{\emergencystretch}{3em}

\newcommand{\lam}{\lambda}
\newcommand{\best}[1]{\textbf{#1}}
\newcommand{\armctrl}{\textsc{2d-ctrl}}
\newcommand{\armref}{\textsc{sv-ref}}

%------------------------------------------------------------------------------
%  Placeholder machinery
%------------------------------------------------------------------------------
\newif\iffinal
\finalfalse   % \finaltrue once every placeholder has been filled
\iffinal
  \newcommand{\PH}[1]{#1}
\else
  \newcommand{\PH}[1]{\textcolor{red}{[#1]}}
\fi
\newcommand{\TBD}{\PH{--}}              % a single missing number
\newcommand{\TODO}[1]{\PH{\textbf{TODO:} #1}}

% \placeholderfig{file}{width}{what the figure should show}
% Includes the file if it exists, otherwise draws a framed box.
\newcommand{\placeholderfig}[3]{%
  \IfFileExists{#1}{%
    \includegraphics[width=#2]{#1}%
  }{%
    \fbox{\parbox[c][0.28\textheight][c]{\dimexpr#2-2\fboxsep-2\fboxrule}{%
      \centering\small\textcolor{red}{\textbf{Figure placeholder}}\\[0.4em]
      \texttt{\detokenize{#1}}\\[0.6em]
      \parbox{0.8\linewidth}{\centering\footnotesize #3}}}%
  }%
}

\title{\textbf{Effect of an Authored Joint-Limit Prior on\\
Single-View 3D Pose Regression Without 3D Supervision}\\[0.4em]
\large A from-scratch, 2D-only study on the SMILySTICKS benchmark}

\author{%
  Youness ANOUAR\\
  \small Forschungszentrum J\"ulich --- SMILify project
}

\date{\PH{date}}

%==============================================================================
\begin{document}
\maketitle

\iffinal\else
\begin{center}
\fcolorbox{red}{red!5}{\parbox{0.9\textwidth}{\small\raggedright
\textbf{Draft placeholder.} Red text marks values and statements still to be
filled from \texttt{prior\_study\_results\_2d/}. The structure, method and
caveats reflect the study design in \texttt{cmnds\_2d\_only.txt},
\texttt{configs\_runs/singleview\_2d\_lam*.json} and
\texttt{docs/ISSUE\_2d\_only\_depth\_flip.md}.}}
\end{center}
\fi

\begin{abstract}
\PH{Two to four sentences: the question, the three-arm design, the headline
plausibility result, the headline N-MPJPE / 2D accuracy result, and the
recommended $\lam$ (if any) for 2D-only training.}
\end{abstract}

\noindent\rule{\textwidth}{0.4pt}

%==============================================================================
\section{Introduction}
\label{sec:intro}

A previous sweep showed that, in 3D-supervised continuation training, an
authored per-joint rotation prior reduces anatomically impossible poses by up to
three orders of magnitude at a small but systematic accuracy cost. That study
could not isolate the value of the prior itself: its arms were continuations of a
3D-supervised checkpoint, trained with a \texttt{keypoint\_3d} term at weight
$20$ that already pinned the pose, and it lacked an epoch-matched $\lam = 0$
control.

When 3D supervision is removed, the 2D keypoint loss is the only data term, the
problem is underdetermined, and the prior becomes one of the few constraints on
the articulation. The study reported here asks:

\begin{quote}
\emph{Without 3D supervision, does an authored per-joint rotation prior produce
more anatomically plausible poses, and does it improve (or preserve) 3D and 2D
accuracy relative to an identical 2D-only model without the prior?}
\end{quote}

\noindent This report gives:
\begin{enumerate}[nosep,leftmargin=1.6em]
  \item the plausibility effect of the prior against an epoch-matched control
        (\Cref{sec:res-plaus});
  \item its effect on scale-corrected 3D accuracy and 2D reprojection accuracy
        (\Cref{sec:res-acc});
  \item a range-of-motion analysis (\Cref{sec:res-rom});
  \item the cost of dropping 3D supervision, measured against the 3D-supervised
        reference (\Cref{sec:res-vs-ref}).
\end{enumerate}

%==============================================================================
\section{Method}
\label{sec:method}

\subsection{Model and data}

All arms use the single-view SMILify regressor: a
\texttt{vit\_large\_\allowbreak patch16\_\allowbreak 224} backbone with a transformer-decoder head
(depth $6$, $8$ heads), a 6D rotation representation, separate
scale/translation prediction and log-parameterised mesh scaling. The dataset is
\texttt{SMILySTICKS\_\allowbreak centred\_\allowbreak reprojected\_\allowbreak FIXED.h5} in camera-centric frame
convention, with all views expanded and \texttt{dataset\_fraction} $= 0.2$. The
split uses ratios $0.85/0.05/0.10$ and seed $42$; the pre-flight check
(\texttt{preflight\_2d\_only.py}) confirms the test split matches that of the
3D-supervised reference, so all arms are scored on the same frames.
\TODO{confirm the number of test samples / view-items actually scored.}

\subsection{Removing 3D supervision}

Each arm is trained \emph{from scratch} with the training setup that produced
the 3D-supervised reference, with exactly two changes: all 3D supervision
weights are set to $0$, and the joint-limit weight is the swept $\lam$.

\begin{table}[t]
\centering
\small
\caption{Loss terms in the 2D-only arms (base weights at epoch 0). Regularizer
weights then follow the reference curriculum unchanged.}
\label{tab:losses}
\resizebox{\textwidth}{!}{%
\begin{tabular}{lll}
\toprule
Status & Terms & Weight \\
\midrule
Dropped (3D)  & \texttt{keypoint\_3d}, \texttt{global\_rot}, \texttt{joint\_rot},
                \texttt{betas}, \texttt{trans}, & $0$ \\
              & \texttt{log\_beta\_scales}, \texttt{betas\_trans} & \\
Data term     & \texttt{keypoint\_2d} & $0.1$ \\
Camera terms  & \texttt{fov}, \texttt{cam\_rot}, \texttt{cam\_trans} & \PH{$0$ in config} \\
Regularizers  & \texttt{joint\_angle}, \texttt{limb\_scale}, \texttt{limb\_trans}
              & $10^{-3}$, $0.1$, $1$ (curriculum) \\
Prior         & \texttt{joint\_limit\_regularization} & $\lam$ \\
\bottomrule
\end{tabular}}
\end{table}

\paragraph{Positive depth.} A first launch of these runs converged to a
behind-the-camera solution: the pinhole projection maps $(x,y,z)$ and
$(-x,-y,-z)$ to the same pixel, so the 2D loss cannot fix the sign of depth.
The runs were restarted with depth parameterised as
$z = z_0 \exp(\hat{z})$, $z_0 = 0.25$, which guarantees $z > 0$ without adding
supervision, and with the mesh scale initialised at $0.04$ to match the
reference (\texttt{fov}, \texttt{cam\_rot} and \texttt{cam\_trans} are also
$0$ in the configs; \TODO{confirm this is intended, since the study notes list
them as kept}). The flipped runs are excluded from this report.
\TODO{confirm with \texttt{diagnose\_mesh\_render.py} that all three restarted
arms have depth $>0$ and no vertices behind the camera.}

\subsection{Joint-limit regularisation}

Let $\theta_{a}$ be the predicted Euler angle on rotation axis $a$ and
$[\,m_a, M_a\,]$ its authored range from
\texttt{3D\_model\_prep/SMILy\_STICK\_limits\_authored.pkl} ($129$ of $162$
non-root axes bounded). The penalty is the mean one-sided hinge
\begin{equation}
  \mathcal{L}_{\text{limit}}
  = \lam \cdot \frac{1}{|\mathcal{A}|}\sum_{a \in \mathcal{A}}
  \Big[\operatorname{relu}(\theta_a - M_a) + \operatorname{relu}(m_a - \theta_a)\Big].
  \label{eq:penalty}
\end{equation}
Because the total loss no longer contains the \texttt{keypoint\_3d} term, a given
$\lam$ is a much larger fraction of the objective than in the earlier sweep;
$\lam$ values are \emph{not} comparable across the two studies.

\subsection{Experimental arms}

\begin{table}[t]
\centering
\small
\caption{Arms of the 2D-only study. All three are trained from scratch with an
identical schedule, allocation and seed, and differ only in $\lam$. The
3D-supervised reference is listed for context only.}
\label{tab:arms}
\begin{tabular}{llcccl}
\toprule
Arm & Supervision & $\lam$ & Epochs & Best epoch & Role \\
\midrule
\armctrl{}       & 2D only & $0$       & \PH{N} & \TBD & control / baseline \\
\textsc{2d-1e-4} & 2D only & $10^{-4}$ & \PH{N} & \TBD & weak prior \\
\textsc{2d-1e-1} & 2D only & $10^{-1}$ & \PH{N} & \TBD & strong prior \\
\midrule
\armref{}        & 2D + 3D & $0$       & 386    & 386  & context only \\
\bottomrule
\end{tabular}
\end{table}

Training ran on RWTH CLAIX-2023 with $2$ nodes $\times$ $4$ H100 ($8$ ranks) per
arm, \texttt{batch\_size} \PH{b} per rank (effective batch \PH{B}), chained
across 24\,h jobs. \TODO{state total epochs and wall-clock per arm; state whether
the effective batch matches the reference (Caveat 5 in
\texttt{cmnds\_2d\_only.txt}).}

\subsection{Metrics}
\label{sec:metrics}

\paragraph{Anatomical plausibility.} Violation rate and overshoot per authored
axis, as in the previous study. These are angle-based and unaffected by global
scale drift.

\paragraph{3D accuracy.} Without 3D supervision nothing anchors absolute size:
a larger insect farther away reprojects to the same pixels. We therefore report
\begin{itemize}[nosep,leftmargin=1.6em]
  \item the \emph{fitted global scale} (predicted / true size; $1.0$ = no drift),
        read first as a diagnostic;
  \item \textbf{N-MPJPE} (one global scale fitted over the test set) ---
        \emph{the headline accuracy metric};
  \item PA-MPJPE (per-frame Procrustes alignment), an optimistic bound on pose
        quality;
  \item raw MPJPE, only alongside the fitted scale.
\end{itemize}

\paragraph{2D accuracy.} PCK curves and mean 2D error at native ($1530$\,px) and
input ($224$\,px) resolution. These are unaffected by the depth and scale
ambiguities.

\paragraph{Range of motion.} Per-axis observed ROM and the clipping bound
$\mathrm{ROM}^{\text{clip}}_a$, computed here against the \emph{control} arm.

%==============================================================================
\section{Results}
\label{sec:results}

\begin{table}[t]
\centering
\small
\setlength{\tabcolsep}{3.8pt}
\caption{Main results on the test split. Deltas in the text are relative to the
control \armctrl{}. \armref{} is 3D-supervised and shown for context.
\PH{Bold the best 2D-only value per column.}}
\label{tab:main}
\resizebox{\textwidth}{!}{%
\begin{tabular}{l c *{10}{r}}
\toprule
 & & \multicolumn{3}{c}{Plausibility} & \multicolumn{4}{c}{3D accuracy} & \multicolumn{3}{c}{2D accuracy} \\
\cmidrule(lr){3-5}\cmidrule(lr){6-9}\cmidrule(lr){10-12}
Arm & $\lam$ &
rate & mean ov. & max ov. &
scale & N-MPJPE & PA-MPJPE & MPJPE &
PCK@5 & PCK@5 & 2D err. \\
 & &
(\%) & (deg) & (deg) &
 & (mm) & (mm) & (mm) &
nat. & 224 & (px) \\
\midrule
\armctrl{}       & 0         & \TBD & \TBD & \TBD & \TBD & \TBD & \TBD & \TBD & \TBD & \TBD & \TBD \\
\textsc{2d-1e-4} & $10^{-4}$ & \TBD & \TBD & \TBD & \TBD & \TBD & \TBD & \TBD & \TBD & \TBD & \TBD \\
\textsc{2d-1e-1} & $10^{-1}$ & \TBD & \TBD & \TBD & \TBD & \TBD & \TBD & \TBD & \TBD & \TBD & \TBD \\
\midrule
\armref{}        & 0 (3D)    & 14.76 & 2.53 & 77.70 & \TBD & \TBD & \TBD & 1.092 & 0.376 & 0.994 & 7.827 \\
\bottomrule
\end{tabular}}
\end{table}

\begin{figure}[t]
\centering
\placeholderfig{fig_2d_tradeoff.pdf}{\textwidth}{Violation rate (log) and
N-MPJPE vs.\ $\lam$ for the three 2D-only arms, with the 3D-supervised
reference as an open marker.}
\caption{Plausibility/accuracy trade-off without 3D supervision.
\PH{Describe the shape of both curves.}}
\label{fig:tradeoff}
\end{figure}

\subsection{Anatomical plausibility}
\label{sec:res-plaus}

\PH{Violation rate falls from X\% (control) to Y\% ($10^{-4}$) and Z\%
($10^{-1}$); maximum overshoot from A to B to C degrees. Compare the
control's violation rate with the reference's 14.76\,\%: does removing 3D
supervision make poses less plausible?}

\begin{figure}[t]
\centering
\placeholderfig{fig_2d_worst_axes.pdf}{0.92\textwidth}{The twelve axes most
violated by the control, and their violation rate under each $\lam$ (log).}
\caption{Most-violated axes in the 2D-only setting. \PH{Which joint groups
dominate? Same as the 3D-supervised case (coxa/femur flexion)?}}
\label{fig:worst}
\end{figure}

\subsection{Accuracy}
\label{sec:res-acc}

\paragraph{Scale drift.} \PH{Fitted global scale per arm. If far from 1.0,
state it and rely on N-MPJPE.}

\paragraph{3D accuracy.} \PH{N-MPJPE and PA-MPJPE vs.\ the control. This is the
central question: in the 2D-only regime, the prior may \emph{improve} 3D
accuracy by resolving depth ambiguities, unlike in the 3D-supervised sweep.
Check that the ordering also holds at the last checkpoint of each arm
(\Cref{sec:disc-validity}).}

\paragraph{2D accuracy.} \PH{PCK and mean 2D error. A drop here means the prior
is overriding the only data term.}

\begin{figure}[t]
\centering
\placeholderfig{fig_2d_pck.pdf}{\textwidth}{PCK curves at native and input
resolution for the three 2D-only arms and the reference.}
\caption{PCK curves on the test split. \PH{Ordering of arms.}}
\label{fig:pck}
\end{figure}

\begin{figure}[t]
\centering
\placeholderfig{fig_2d_nmpjpe_quantiles.pdf}{\textwidth}{N-MPJPE quantile
functions (left) and change relative to the control per percentile (right).}
\caption{N-MPJPE by percentile. \PH{Does the prior's effect concentrate in the
tail, where implausible poses live?}}
\label{fig:quant}
\end{figure}

\subsection{Range of motion}
\label{sec:res-rom}

\begin{table}[t]
\centering
\small
\caption{Mean range of motion per bounded axis (deg), by joint group, over the
$129$ bounded axes. The clipping bound is computed from the control arm.}
\label{tab:rom}
\begin{tabular}{l r *{4}{r}}
\toprule
Joint group & axes & \armctrl{} & $10^{-4}$ & $10^{-1}$ & \armref{} \\
\midrule
Coxa            & 18 & \TBD & \TBD & \TBD & 78.1 \\
Trochanter      & 18 & \TBD & \TBD & \TBD & 74.2 \\
Femur           & 18 & \TBD & \TBD & \TBD & 71.6 \\
Tibia           & 18 & \TBD & \TBD & \TBD & 88.9 \\
Tarsus          & 18 & \TBD & \TBD & \TBD & 96.3 \\
Antenna         & 18 & \TBD & \TBD & \TBD & 40.9 \\
Head / wings    & 21 & \TBD & \TBD & \TBD & 20.2 \\
\midrule
\textbf{All}    & \textbf{129} & \TBD & \TBD & \TBD & \textbf{66.1} \\
\midrule
\multicolumn{2}{l}{\emph{Clipping bound (control)}} & \multicolumn{4}{c}{\TBD} \\
\multicolumn{2}{l}{\emph{Mean authored range}} & \multicolumn{4}{c}{$51.6$} \\
\bottomrule
\end{tabular}
\end{table}

\PH{Does the 2D-only control over- or under-articulate relative to the
reference? Does the prior again induce a learned inward margin below the
clipping bound (femur/tibia)?}

\begin{figure}[t]
\centering
\placeholderfig{fig_2d_rom_retention.pdf}{0.86\textwidth}{ROM retained
relative to the control, grouped by the control's violation rate, with the
clipping bound per group.}
\caption{ROM retention vs.\ clipping bound. \PH{Summary.}}
\label{fig:retention}
\end{figure}

\begin{figure}[t]
\centering
\placeholderfig{fig_2d_rom_ratio.pdf}{0.86\textwidth}{Distribution over the 129
bounded axes of observed ROM / authored range, per arm.}
\caption{Use of the authored envelope. \PH{Mean/median ratio per arm; number of
axes above 1.0.}}
\label{fig:ratio}
\end{figure}

\subsection{Cost of removing 3D supervision}
\label{sec:res-vs-ref}

\PH{Control vs.\ \armref{}: plausibility, N-MPJPE, PA-MPJPE, PCK. This
difference measures ``no 3D supervision'' only. Does the best constrained arm
close part of the gap?}

\subsection{Qualitative results}
\label{sec:res-qual}

\begin{figure}[t]
\centering
\placeholderfig{fig_2d_render_qualitative.png}{\textwidth}{Rendered exported
poses of all three arms on the same worst-case windows (one per camera),
selected from the control arm's export.}
\caption{Qualitative comparison on identical frames. \PH{What is visible:
implausible legs in the control? Stiff gait at $\lam=10^{-1}$? Correct depth
ordering?}}
\label{fig:qual}
\end{figure}

%==============================================================================
\section{Discussion}
\label{sec:discussion}

\subsection{Does the prior help when 3D supervision is absent?}
\PH{Interpretation of \Cref{sec:res-plaus,sec:res-acc}, contrasted with the
3D-supervised sweep, where the accuracy cost grew monotonically with $\lam$.}

\subsection{Over-regularisation}
\PH{Is the learned-margin effect stronger here, given that $\lam$ is a larger
share of the loss?}

\subsection{Choice of operating point}
\PH{Recommended $\lam$ for 2D-only training, or ``further sweep needed''
(only two non-zero values were tested).}

\subsection{Threats to validity}
\label{sec:disc-validity}
\begin{itemize}[leftmargin=1.6em,itemsep=0.25em]
  \item \textbf{Scale/depth ambiguity.} Positive depth fixes only the sign; the
        size--distance ambiguity remains. Arms are compared on N-MPJPE.
  \item \textbf{Checkpoint selection.} \texttt{best\_model.pth} is chosen on
        total validation loss, which includes the limit penalty, so constrained
        and control arms are selected under different objectives.
        \PH{Report whether conclusions hold at the last checkpoint.}
  \item \textbf{Non-comparable $\lam$.} The $\lam$ values do not correspond to
        those of the 3D-supervised sweep; the two tables should not be merged.
  \item \textbf{Single seed.} Each arm was trained once (seed 42).
        \PH{State whether training dynamics differed between arms.}
  \item \textbf{Optimisation regime.} \PH{Effective batch vs.\ reference.}
\end{itemize}

%==============================================================================
\section{Conclusion}
\label{sec:conclusion}

\begin{description}[leftmargin=1.4em,style=nextline]
\item[More anatomically plausible poses without 3D supervision --- \PH{yes/no}.]
  \PH{One or two sentences with the key numbers.}
\item[Accuracy relative to the 2D-only control --- \PH{improved/preserved/lost}.]
  \PH{One or two sentences with N-MPJPE and PCK.}
\end{description}

\PH{Closing paragraph: what the prior is worth in the 2D-only regime, and how
this compares with the 3D-supervised case.}

%==============================================================================
\section{Next steps}
\label{sec:next}

\begin{enumerate}[leftmargin=1.9em,itemsep=0.45em]
\item \PH{Denser $\lam$ sweep around the best value.}
\item \PH{Repeat with additional seeds.}
\item \PH{Remove the scale ambiguity (\texttt{allow\_mesh\_scaling = false}) or
      add a size prior.}
\item \PH{Multi-view 2D-only variant.}
\end{enumerate}

%==============================================================================
\appendix
\section{Data provenance}
\label{app:prov}

\begin{description}[leftmargin=0em,style=nextline,itemsep=0.25em,font=\normalfont\ttfamily\small]
\item[configs\_runs/singleview\_2d\_lam\{0,1e-4,1e-1\}.json]
  Training configurations of the three arms.
\item[cmnds\_2d\_only.txt, hpc\_files/rwth/run\_2d\_only\_train.sbatch]
  Study design, launch and scoring procedure.
\item[docs/ISSUE\_2d\_only\_depth\_flip.md]
  Diagnosis of the behind-camera solution and the positive-depth fix.
\item[prior\_study\_results\_2d/sweep/\{sweep.md, sweep.csv\}]
  Aggregated metrics; source of \Cref{tab:main}.
\item[prior\_study\_results\_2d/lam\{1e-4,1e-1\}/comparison/comparison.md]
  Control-vs-$\lam$ comparisons (regenerated with \texttt{--from-scratch}).
\item[prior\_study\_results\_2d/\{sv\_reference, lam*/sv\_constrained\}/analysis/]
  Per-arm violation, per-axis and magnitude statistics; source of the ROM
  analyses.
\item[prior\_study\_results\_2d/renders/segments.json]
  Windows used for \Cref{fig:qual}, picked from the control arm's export.
\item[3D\_model\_prep/SMILy\_STICK\_limits\_authored.pkl]
  Authored per-joint rotation limits.
\end{description}

\noindent\TODO{add the figure-generation script once written.}

\end{document}
